\chapter{MATHEMATICAL MODELING AND METHODS}

\section{Introduction}
The proposed ANFIS-LLFSR framework is a multi-stage hybrid system designed to address the challenges of face restoration under severe low-light conditions. The architecture integrates statistical feature extraction, fuzzy logic-based decision making, and deep learning-based spatial refinement. This chapter details the mathematical formulations and architectural components of the system.

\subsection{The Evolution of the ANFIS-Swin Hybrid}
The design of the ANFIS-LLFSR was not a single-step process but an evolutionary refinement. Initial prototypes utilized a standard Residual-in-Residual Dense Block (RRDB) backbone. However, we observed that RRDBs, while effective for natural textures, failed to maintain the long-range structural symmetry of faces (e.g., matching the eye-geometry across a shadowed midline). This led to the integration of the Swin Transformer, which provides the necessary global receptive field. The ANFIS meta-controller was retained from the early design to provide the interpretable "illumination logic" that pure transformers often lack when trained on diverse datasets.

\subsection{Rejected Architectural Alternatives and Lessons Learned}
During the development phase, several alternative configurations were evaluated and rejected:
\begin{enumerate}
    \item \textbf{Pure GAN-based Enhancement:} While visually sharper, pure GANs exhibited unacceptable "Identity Drift" in severe L3 cases ($DF > 0.8$), occasionally hallucinating features from the training set that did not belong to the subject.
    \item \textbf{Heuristic Gamma Correction:} Pre-processing with a fixed gamma curve was found to amplify sensor noise in deep shadows while washing out detail in ambient regions.
    \item \textbf{Non-Local Networks:} These were rejected due to their $\mathcal{O}(N^2)$ complexity, which made $512 \times 512$ restoration infeasible on consumer hardware. The window-based Swin approach was selected as the optimal trade-off between global modeling and local efficiency.
\end{enumerate}

\begin{figure}[htbp]
    \centering
    \includegraphics[width=0.9\textwidth]{thesis/figures/system_flowchart}
    \caption{End-to-End System Architecture of ANFIS-LLFSR, showing the interaction between the ANFIS Meta-Controller and the SwinFuzzyLCR Restoration Layer.}
    \label{fig:system_flow}
\end{figure}

\section{Data Preprocessing and Feature Extraction}
\subsection{Statistical Feature Extraction for ANFIS}
The effectiveness of the ANFIS module depends on the quality of the input features. We extract a 10-dimensional feature vector $\mathbf{x}_{feat}$ that captures both the spatial and frequency-domain characteristics of the low-light degradation:

\begin{enumerate}
    \item \textbf{Global Mean Intensity ($\mu_I$):} The average pixel value across all channels, representing the baseline illumination.
    \item \textbf{Local Contrast ($\sigma_I$):} The standard deviation of pixel intensities, indicating the dynamic range.
    \item \textbf{Dark Pixel Ratio ($R_{dark}$):} The fraction of pixels with values below a threshold (e.g., 0.1), highlighting the severity of clipping in shadows.
    \item \textbf{Entropy ($H$):} A measure of the information content and noise level in the image.
    \item \textbf{High-Frequency Energy ($E_{HF}$):} Calculated from the Sobel-gradient magnitude, indicating the presence of structural details.
    \item \textbf{Color Saturation Mean:} The average saturation in HSV space, which often drops in low-light captures.
    \item \textbf{Blur Metric ($B$):} Based on the Laplacian variance, estimating the degree of motion or defocus blur.
    \item \textbf{Noise Variance ($\sigma^2_n$):} Estimated using a median absolute deviation (MAD) approach in the wavelet domain.
    \item \textbf{Histogram Skewness:} Measuring the asymmetry of the intensity distribution.
    \item \textbf{Spatial Coherence:} A local binary pattern (LBP) based feature that captures the micro-structural integrity.
\end{enumerate}

These features are normalized to the range $[0, 1]$ before being passed to the ANFIS fuzzification layer.

\subsection{Degradation Analysis and Manifold Visualization}
To understand the distribution of these features, we performed a t-SNE (t-Distributed Stochastic Neighbor Embedding) visualization of the 10-D feature manifold. This revealed that images with similar darkness factors cluster together, but with distinct sub-clusters based on the type of noise (Poisson vs. Gaussian). This clustering behavior justifies the use of a fuzzy inference system, as it can define membership functions that "wrap" around these clusters in feature space, providing a smooth transition between different degradation states.

\subsection{Synthetic Degradation Pipeline}
To train a robust model, we simulate realistic low-light degradations using a compound pipeline that mimics the physical characteristics of low-light imaging:
\begin{enumerate}
    \item \textbf{Resolution Reduction:} The HR image $I$ is downsampled by a factor $s \in \{4, 8\}$ to create the LR base $I_{LR}$.
    \item \textbf{Illumination Transformation:} We apply a non-linear mapping $f(I) = I^\gamma$, where $\gamma$ is sampled from a Gaussian distribution $\mathcal{N}(3.5, 0.5)$.
    \item \textbf{Shot Noise Modeling:} Poisson noise is added to simulate photon arrival statistics, followed by Gaussian read noise.
    \item \textbf{Motion Blur Kernel:} A random motion blur kernel is generated using a 2D random walk process, simulating camera shake.
\end{enumerate}

\section{Baseline Model Selection and Rationale}
In alignment with standard academic practices, we selected a set of baseline models that represent the historical progression of the field.

\subsection{SRCNN: Convolutional Super-Resolution}
SRCNN is used as the simplest deep learning baseline. It demonstrates the basic mapping from LR to HR spaces. We use it to quantify the performance gain provided by the addition of fuzzy logic and transformers.

\subsection{Zero-DCE: Illumination Benchmarking}
Zero-DCE is selected to validate the "enhancement" part of our pipeline. By comparing our ANFIS-guided enhancement with Zero-DCE's curve-based approach, we can isolate the contribution of the fuzzy logic module in managing brightness uncertainty.

\subsection{CodeFormer: Structural SOTA}
CodeFormer serves as the state-of-the-art benchmark for structural face restoration. We use it to evaluate whether our LCR-based hallucination can match or exceed the performance of discrete codebook lookups in terms of identity preservation.

\section{The Neuro-Fuzzy Meta-Controller (ANFIS)}
The ANFIS module serves as the primary illumination estimator and architectural "Manager." It calculates a continuous "Darkness Factor" ($DF$) that guides the subsequent enhancement and reconstruction stages. Unlike black-box neural networks, the ANFIS module provides an interpretable decision trace through its fuzzy rule base.

\subsection{Fuzzification: Mapping the Darkness Manifold}
In the first layer, input features $\mathbf{x} = [x_1, x_2, \dots, x_n]$ are mapped to fuzzy sets using Gaussian membership functions (MFs). For each input feature $i$ and rule $k$, the membership degree is:
\begin{equation}
    \mu_{i,k}(x_i) = \exp \left( -\frac{(x_i - c_{i,k})^2}{2\sigma_{i,k}^2} \right)
\end{equation}
where $\{c_{i,k}, \sigma_{i,k}\}$ are the premise parameters representing the center and spread of the $k$-th fuzzy set for the $i$-th input feature. 

\textbf{Variable Interpretation and Parameter Sensitivity:} The parameter $c_{i,k}$ determines the "focal point" of the illumination regime (e.g., a center at 0.1 for the "Dark" set). The spread parameter $\sigma_{i,k}$ controls the "softness" of the transition; a high $\sigma$ implies a broad illumination category that overlaps significantly with its neighbors, while a low $\sigma$ represents a specialized rule that fires only for a narrow range of intensity features. During training, these parameters are optimized via gradient descent, allowing the MFs to "drift" and "stretch" until they perfectly partition the 10-D feature manifold.
    
\subsection{Layer 2: Rule Firing Strength}
Each rule $r$ in the rule base represents a specific combination of input membership degrees. The firing strength $w_r$ is calculated using the product T-norm:
\begin{equation}
    w_r = \prod_{i=1}^n \mu_{i, r_i}(x_i)
    \label{eq:firing}
\end{equation}
where $r_i$ denotes the index of the MF used for input $i$ in rule $r$.

\textbf{Engineering Intuition:} Equation \ref{eq:firing} represents the "Logical AND" operation. It ensures that a rule only fires if *all* its constituent conditions are met. For example, a rule for "Extreme Low Light" will only have a high $w_r$ if the Mean Intensity is low AND the Dark Pixel Ratio is high. This multi-conditional triggering prevents the system from making aggressive illumination adjustments based on a single noisy feature.

\subsection{Layer 3: Normalization}
The firing strengths are normalized to ensure they sum to unity:
\begin{equation}
    \bar{w}_r = \frac{w_r}{\sum_{j=1}^R w_j}
    \label{eq:norm}
\end{equation}
where $R$ is the total number of rules.

\textbf{Implementation Rationale:} Normalization (Eq \ref{eq:norm}) is critical for the stability of the meta-controller. It ensures that the total "influence" of the fuzzy rules is always 1.0, regardless of the absolute values of the firing strengths. This acts as a soft-max layer that provides a relative weight to each rule, enabling the "Fuzzy Blending" of different illumination enhancement strategies.

\subsection{Layer 4: Consequent Layer}
We utilize a Takagi-Sugeno type fuzzy system, where each rule $r$ has a linear consequent function:
\begin{equation}
    f_r(\mathbf{x}) = p_{r,0} + \sum_{i=1}^n p_{r,i} x_i
\end{equation}
where $\{p_{r,i}\}$ are the consequent parameters.

\subsection{Defuzzification and Rule Visualization}
The final crisp output, the Darkness Factor $DF$, is the weighted sum of all rule outputs:
\begin{equation}
    DF = \sum_{r=1}^R \bar{w}_r f_r(\mathbf{x})
\end{equation}

\begin{figure}[htbp]
    \centering
    \includegraphics[width=0.8\textwidth]{thesis/figures/anfis_surface}
    \caption{ANFIS decision surface visualizing the non-linear mapping between Global Intensity ($\mu_I$) and Local Contrast ($\sigma_I$) to the Darkness Factor ($DF$).}
    \label{fig:anfis_surface}
\end{figure}

The ANFIS module is trained using a hybrid algorithm: Recursive Least Squares (RLS) for consequent parameters and Gradient Descent for premise parameters.

\subsection{Premise Parameter Optimization via Gradient Descent}
The premise parameters $\{c_{i,k}, \sigma_{i,k}\}$ in Layer 1 define the shape of the fuzzy sets. During the backward pass, the error $E = \frac{1}{2}(y - y_{target})^2$ is backpropagated using the chain rule:
\begin{equation}
    \frac{\partial E}{\partial c_{i,k}} = \frac{\partial E}{\partial y} \frac{\partial y}{\partial \bar{w}_r} \frac{\partial \bar{w}_r}{\partial w_r} \frac{\partial w_r}{\partial \mu_{i,k}} \frac{\partial \mu_{i,k}}{\partial c_{i,k}}
\end{equation}
This allows the membership functions to adaptively shift and stretch to better cover the input feature space.

\subsection{Consequent Parameter Optimization via Recursive Least Squares (RLS)}
To achieve rapid convergence, the consequent parameters are optimized analytically using RLS. Given the normalized firing strengths $\bar{w}$ and the augmented input vector $\mathbf{x}_{aug} = [\mathbf{x}, 1]$, the output can be written as:
\begin{equation}
    y = \sum_{r=1}^R (\bar{w}_r \mathbf{x}_{aug}^T) \mathbf{p}_r
\end{equation}
By concatenating the parameters into a global vector $\mathbf{P}$, we solve the linear problem $\mathbf{y} = \mathbf{A} \mathbf{P}$. The RLS algorithm updates the estimate $\mathbf{P}_t$ as:
\begin{equation}
    \mathbf{P}_t = \mathbf{P}_{t-1} + \mathbf{S}_t \mathbf{a}_t (y_t - \mathbf{a}_t^T \mathbf{P}_{t-1})
\end{equation}
where $\mathbf{S}_t$ is the covariance matrix updated using the matrix inversion lemma. This "hybrid" approach allows the model to reach high accuracy ($>90\%$) in significantly fewer epochs than pure gradient descent.

\section{Structural Anchoring via Locality Constrained Representation (LCR)}
The LCR module reconstructs high-frequency face patches using a coupled LR-HR dictionary.

\subsection{LCR Encoding and Sparse Priors}
Given an LR patch $p$, we find its representation coefficients $\alpha^*$ by solving the locality-constrained optimization problem:
\begin{equation}
    \min_\alpha \|p - D_{LR}\alpha\|^2 + \lambda \|d \odot \alpha\|^2
\end{equation}
where $d$ is the locality weight vector, $d_k = \exp(\|p - D_{LR,k}\|^2 / \beta)$, and $\lambda$ is a regularization parameter. The closed-form solution is:
\begin{equation}
    \alpha^* = (D_{LR}^T D_{LR} + \lambda \text{diag}(d^2))^{-1} D_{LR}^T p
\end{equation}

\subsection{Closed-Form Optimization and Complexity Bottlenecks}
\textbf{Computational Tradeoffs:} The $1024^3$ complexity of the matrix inversion is the primary bottleneck. We mitigate this through \textit{Pre-Computed Dictionary Covariances} ($D_{LR}^T D_{LR}$), reducing the per-patch cost to an $\mathcal{O}(N^2)$ update. This implementation insight allowed us to reduce the inference time by 40\% compared to the naive LCR approach.

\subsection{Theoretical Bounds of Sparse Reconstruction in Low-Light}
In the L3 (Severe) regime, the signal-to-noise ratio (SNR) is often below the theoretical threshold for stable sparse recovery. Our system addresses this through the \textit{Adaptive Locality Bandwidth} ($\beta_{adaptive}$), which dynamically expands the search radius in the dictionary as the Darkness Factor increases. This prevents "patch-collapse" where a dark input is incorrectly matched to a noisy atom, ensuring that the hallucination remains photometrically consistent.
    

\section{System Hyperparameters and Configuration}
The performance of the ANFIS-LLFSR system is highly dependent on the choice of hyperparameters across its multi-stage pipeline. Table \ref{tab:hyperparams_detailed} provides a complete configuration used for our optimal "Research-Grade" model.

\begin{table}[htbp]
    \centering
    \caption{Detailed System Hyperparameters for ANFIS-LLFSR.}
    \label{tab:hyperparams_detailed}
    \begin{tabular}{lll}
        \toprule
        \textbf{Module} & \textbf{Hyperparameter} & \textbf{Value} \\
        \midrule
        \textbf{ANFIS} & Input Features & 10-D Vector \\
                       & Num. Membership Functions & 3 per input (Low, Med, High) \\
                       & MF Type & Gaussian ($\sigma, c$) \\
                       & Learning Rule & Hybrid (RLS + SGD) \\
        \midrule
        \textbf{LCR}   & Dictionary Size ($N_{atoms}$) & 1024 patches (LR-HR pairs) \\
                       & Patch Size & 8x8 pixels \\
                       & Locality Bandwidth ($\beta$) & 0.12 (Base) \\
                       & Stride & 4 pixels \\
        \midrule
        \textbf{SwinIR} & Embedding Dim & 180 \\
                       & Num. Blocks & 6 Residual Swin Transformer Blocks \\
                       & Attention Heads & 6 \\
                       & Window Size & 8x8 \\
        \midrule
        \textbf{Losses} & $\lambda_{pixel}$ (Charbonnier) & 1.0 \\
                       & $\lambda_{percep}$ (VGG-19) & 0.1 \\
                       & $\lambda_{adv}$ (LSGAN) & 0.05 \\
                       & $\lambda_{ID}$ (ArcFace) & 0.08 \\
        \bottomrule
    \end{tabular}
\end{table}

\section{Computational Complexity and Runtime Profile}
To evaluate the feasibility of forensic deployment, we analyzed the computational overhead of the multi-stage pipeline. The inference process is divided into three primary bottlenecks:

\begin{enumerate}
    \item \textbf{LCR Patch Search:} The time complexity of the dictionary search is $\mathcal{O}(N_{patches} \cdot N_{atoms}^2)$ after pre-computing the covariance matrix. While theoretically heavy, the use of patch-parallel processing on the GPU allows for $512 \times 512$ restoration in $\approx 350$ms.
    \item \textbf{Swin Transformer Refinement:} The self-attention mechanism in the Swin blocks has a complexity of $\mathcal{O}(H \cdot W \cdot C^2)$, where $C$ is the embedding dimension. Our optimized configuration (6 blocks, 6 heads) consumes approx. 2.4GB of VRAM during the forward pass.
    \item \textbf{ANFIS Meta-Control:} The neuro-fuzzy inference stage is computationally negligible ($\mathcal{O}(N_{rules} \cdot N_{features})$), consuming less than 1ms per frame.
\end{enumerate}

The total end-to-end inference time is 0.82 seconds on an NVIDIA RTX 3090, making the system suitable for near-real-time forensic review where high-resolution accuracy is prioritized over high-framerate streaming.

\section{Identity Preservation Theory: The Role of ArcFace}
A critical novelty of our framework is the use of ArcFace embeddings to anchor the hallucination process.

\subsection{Additive Angular Margin Loss}
Unlike standard Softmax loss, which only ensures feature separability, ArcFace introduces an additive angular margin to maximize intra-class compactness and inter-class discrepancy. This creates a hyperspherical embedding space where facial identity is represented as a high-dimensional vector.

\subsection{ArcFace as a Geometric Sentinel}
In our pipeline, the ArcFace loss acts as a sentinel. If the LCR hallucination or Swin refinement begins to "drift" toward a generic face, the ArcFace gradient pulls the features back toward the original subject's embedding. This ensures that the restored face is not just "a face" but "the same face" as the input.

\section{Training Procedure and Execution Loop}
The implementation of the training loop follows a progressive curriculum designed to stabilize the gradients across the hybrid architecture. The execution flow in \texttt{train\_full.py} follows three distinct phases:
\begin{enumerate}
    \item \textbf{Phase 1: ANFIS Warm-up:} The illumination estimator is trained independently on a synthetic darkness dataset for 50 epochs to establish a reliable baseline for the $DF$ signal.
    \item \textbf{Phase 2: Joint Optimization:} The SwinFuzzyLCR refinement network is introduced, and the entire pipeline is optimized end-to-end. During this phase, the ANFIS weights are frozen for the first 10,000 iterations to stabilize the transformer's initial gradients.
    \item \textbf{Phase 3: Identity Hardening:} The ArcFace loss is activated, and the learning rate is decayed by a factor of 0.5 every 50,000 iterations using a StepLR scheduler.
\end{enumerate}

\subsection{Addressing Training Divergence in Multi-Loss Optimization}
A major engineering challenge was the "Loss Conflict" between the LSGAN adversarial loss and the ArcFace identity loss. Early in training, the GAN discriminator would force the generator to produce high-frequency noise to "look real," which frequently perturbed the identity embeddings. We solved this by implementing a \textit{Gradient Scaling} strategy: the identity loss gradient was scaled by a factor of 2.0 whenever the adversarial loss began to dominate. This "Identity Over-riding" mechanism ensures that subject fidelity is never sacrificed for superficial sharpness.

\section{Membership Function Design and Sensitivity Analysis}
The choice of membership function (MF) shape significantly impacts the "softness" of the illumination transitions.

\subsection{Gaussian vs. Triangular Membership Functions}
We compared the performance of Gaussian MFs against traditional Triangular MFs. While Triangular MFs are computationally cheaper, Gaussian MFs provide a differentiable surface that is better suited for gradient-based optimization. Furthermore, Gaussian MFs do not "saturate" at 0 or 1 as abruptly, allowing for a more nuanced representation of the 10-D feature manifold.

\subsection{Impact of MF Overlap}
The degree of overlap between adjacent MFs determines the "fuzziness" of the decision boundary. We found that an overlap of 25\% across the feature range $[0, 1]$ prevents "dead zones" where no rules are fired, while maintaining sufficient rule specialization for different illumination regimes.

\section{Recursive Least Squares (RLS) Convergence Properties}
The hybrid learning algorithm utilizes RLS to solve for the consequent parameters in Layer 4. 

\subsection{Analytic Solution vs. Iterative Search}
Unlike stochastic gradient descent (SGD), RLS provides an analytic solution to the least-squares problem at each step. This allows the ANFIS module to "snap" to the optimal linear surface for a given set of premise parameters. Mathematically, the update follows:
\begin{equation}
    P_k = P_{k-1} + \frac{S_{k-1} a_k (y_k - a_k^T P_{k-1})}{1 + a_k^T S_{k-1} a_k}
\end{equation}
where $S_k$ is the covariance matrix. This rapid convergence is essential for the multi-stage training pipeline, as it allows the ANFIS signals to stabilize quickly before the transformer refinement network begins its optimization.

\section{Transformer-based Feature Refinement}
The final stage of the pipeline utilizes the Swin Transformer architecture for spatial refinement.

\subsection{Window-based Self-Attention}
To maintain computational efficiency on high-resolution images ($512 \times 512$), the self-attention is calculated within local windows of size $M \times M$ (typically $8 \times 8$). For a window feature map $X \in \mathbb{R}^{M^2 \times C}$, the attention is:
\begin{equation}
    \text{Attention}(Q, K, V) = \text{SoftMax}\left(\frac{QK^T}{\sqrt{d}} + B\right)V
\end{equation}
where $Q, K, V$ are query, key, and value matrices, and $B$ is the relative position bias.

\subsection{Shifted Window Partitioning}
To enable cross-window communication, the Swin Transformer employs a "shifted window" approach. In Layer $l$, the window partitioning is regular. In Layer $l+1$, the windows are shifted by $(\lfloor M/2 \rfloor, \lfloor M/2 \rfloor)$ pixels. This allows the network to gradually learn global structural constraints through local interactions.

\subsection{SwinIR for Super-Resolution}
Our implementation adapts the SwinIR backbone for the LLFSR task. Unlike general SR, our version includes the SFT modulation layers at each stage of the transformer hierarchy, ensuring that the "fuzzy signal" from ANFIS guides the attention heads. This prevents the transformer from attending to noisy or dark artifacts that might otherwise be misinterpreted as structural features.

\subsection{Adaptive Locality Selection Logic}
A unique feature of our LCR implementation is the adaptive selection of the locality bandwidth $\beta$. In traditional LCR, $\beta$ is a fixed hyperparameter. However, we found that for low-light faces, the optimal $\beta$ depends on the signal-to-noise ratio (SNR).

\subsection{Signal-Dependent Bandwidth Modulation}
We utilize the Darkness Factor $DF$ from the ANFIS module to modulate $\beta$:
\begin{equation}
    \beta_{adaptive} = \beta_{base} \cdot (1 + \eta \cdot DF)
\end{equation}
where $\eta$ is a scaling constant. As the image becomes darker ($DF \to 1$), the search radius increases, allowing the dictionary to consider a broader range of "candidate" atoms to compensate for the higher noise in the LR patch.

\subsection{Weighted Dictionary Sampling}
Furthermore, the ANFIS module identifies the "illumination regime" of the input. We partition the patch dictionary $D$ into sub-dictionaries $\{D_{well}, D_{shadow}, D_{dark}\}$. During reconstruction, the LCR algorithm prioritizes atoms from the sub-dictionary that matches the input's $DF$. This ensures that the hallucinated textures are photometrically consistent with the global illumination state.

\subsection{ANFIS-Guided Reweighting}
The coefficients $\alpha^*$ are adjusted using the $DF$ from the ANFIS module. We define an affinity score between the query $DF$ and the precomputed darkness factor of each dictionary atom $DF_{atom, k}$:
\begin{equation}
    \text{Affinity}_k = 1 - |DF - DF_{atom, k}|
\end{equation}
The final reweighted coefficients are:
\begin{equation}
    \alpha_{rw} = \frac{\alpha^* \odot \text{Affinity}}{\|\alpha^* \odot \text{Affinity}\|_1}
\end{equation}
This ensures that patches are reconstructed using atoms that match the global illumination context.

\section{SwinFuzzyLCR Refinement Network}
The deep refinement stage utilizes a transformer-based architecture to polish the initial LCR reconstruction.

\subsection{Spatial Feature Transform (SFT)}
The SFT layer acts as the "Fuzzy Router." It modulates the intermediate feature maps $F$ using the condition vector $c$ (which includes the ANFIS output):
\begin{equation}
    \text{SFT}(F, c) = F \odot \gamma(c) + \beta(c)
\end{equation}
where $\gamma$ and $\beta$ are learned scaling and shifting functions.

\subsection{Convolutional Block Attention Module (CBAM)}
To focus on sensitive facial structures, we apply CBAM, which consists of Channel Attention and Spatial Attention. This allows the model to selectively emphasize feature maps that contain critical identity markers:
\begin{equation}
    F_{refined} = M_s(M_c(F) \otimes F) \otimes (M_c(F) \otimes F)
\end{equation}
where $M_c$ is the channel attention map and $M_s$ is the spatial attention map. 

\subsection{Wavelet-Frequency Fusion for Texture Recovery}
Low-light and blurry images suffer from a collapse of high-frequency components in the Fourier domain. To counteract this, we incorporate a Wavelet-Frequency Fusion block. The block performs a learnable decomposition of features into low-pass (smooth) and high-pass (edge) bands:
\begin{equation}
    F_{LF} = \text{AvgPool}(F), \quad F_{HF} = F - F_{LF}
\end{equation}
The high-frequency features are then refined using a dedicated residual branch before being fused back:
\begin{equation}
    F_{fused} = \text{Fusion}(\text{Concat}(\phi_{LF}(F_{LF}), \phi_{HF}(F_{HF}))) + F
\end{equation}

\textbf{Engineering Rationale for Wavelet Decomposition:} The decision to use a wavelet-based fusion instead of standard residual learning is driven by the "spectral gap" in low-light imagery. In dark captures, the high-frequency components (edges, textures) are often buried under Poisson noise. By explicitly separating the $F_{HF}$ band, we can apply the CBAM attention module specifically to the structural components without being distracted by the low-frequency illumination "glow" in the $F_{LF}$ band. This spectral decoupling prevents the "ringing" artifacts common in global enhancement networks.
    This architecture encourages the model to explicitly reason about facial pores, wrinkles, and hair textures that are often lost during bicubic upsampling.

\section{Identity Preservation via ArcFace Sentinel}
A unique contribution of this work is the use of an "ArcFace Sentinel" to anchor the subject's identity. 

\subsection{ArcFace Embedding Space}
ArcFace utilizes an Additive Angular Margin loss to learn highly discriminative facial embeddings $e \in \mathbb{R}^{512}$. We project these embeddings into the generator's feature space to provide a "subject-specific" bias:
\begin{equation}
    F_{id} = \text{MLP}(e_{ArcFace})
\end{equation}
This vector is added to the deep features, effectively acting as a global prior that tells the network \textit{who} it is reconstructing.

\subsection{Identity Loss Formulation}
The identity loss $\mathcal{L}_{ID}$ is computed as the cosine distance between the embeddings of the reconstructed image $\hat{I}$ and the ground-truth image $I$:
\begin{equation}
    \mathcal{L}_{ID} = 1 - \frac{\text{Enc}(\hat{I}) \cdot \text{Enc}(I)}{\|\text{Enc}(\hat{I})\| \|\text{Enc}(I)\|}
\end{equation}
By minimizing this loss, we ensure that the restoration process does not perturb the high-dimensional identity manifold.

\section{Loss Functions}
The network is trained using a composite loss function to ensure fidelity, perceptual quality, and identity preservation:
\begin{equation}
    \mathcal{L}_{total} = \lambda_1 \mathcal{L}_{pixel} + \lambda_2 \mathcal{L}_{perceptual} + \lambda_3 \mathcal{L}_{adv} + \lambda_4 \mathcal{L}_{ID}
\end{equation}
\begin{itemize}
    \item $\mathcal{L}_{pixel}$: Charbonnier loss for pixel-wise accuracy.
    \item $\mathcal{L}_{perceptual}$: VGG-based perceptual loss.
    \item $\mathcal{L}_{ID}$: Cosine similarity loss between ArcFace embeddings of the reconstructed and ground-truth face.
\end{itemize}

\section{Integrated Inference Workflow and Information Flow}
The ANFIS-LLFSR architecture operates through a "Signal-Modulation-Refinement" loop. Unlike traditional end-to-end black-boxes, our system maintains a clear hierarchical separation between global context and local texture.

\subsection{Hierarchical Information Flow}
The process begins with the **Data Ingestion Tier**, which extracts the 10-D illumination vector. This signal is passed to the **ANFIS Meta-Controller (Manager)**, which generates the Darkness Factor ($DF$). The $DF$ does not modify the image directly; instead, it provides a "Modulation Field" that steers the **SwinFuzzyLCR Worker (Tier 2)**. 

\subsection{Implementation Rationale: Preventing the "Average Face" Problem}
This design choice specifically addresses the "Average Face" limitation. In models trained on datasets with mixed lighting, the weights often converge to a blurry mean that is suboptimal for both dark and bright inputs. By using the SFT layers to "route" the feature maps based on the fuzzy illumination signal, we allow the network to specialize its response. In dark regions ($DF \to 1$), the attention heads are biased toward structural anchoring, while in ambient regions ($DF \to 0$), they focus on high-frequency textural recovery. This ensures that the restoration is context-aware and identity-preserving across the entire illumination manifold.
